\printedpage{49}
% head: A. Grothendieck : Dérivateurs, Ch. I, Édition des Universités de Montpellier II et Paris VII
% foot: Version 8 juillet 2001

To prove that $f'$ is a cohomological $\DD$-fibration, we look over $Y'$ at the functors of the form $\{y'\} \to Y'/y'$ ($y' \in \Ob\, Y'$) and $Y'/y'_1 \to Y'/y'_0$ (if $u' : y'_1 \to y'_0$ an arrow of $Y'$), and we see that the arrows in $\Cat/Y'$ deduced from these by the base change $X' \xrightarrow{f'} Y'$ remain $\DD$-equivalences

\msp{page 80}

\[
\begin{tikzcd}[column sep=large, row sep=large]
X'_1 \arrow[r, "q"] \arrow[d, "f'_1"'] \arrow[dr, phantom, "\text{cart}"] & X'_0 \arrow[r] \arrow[d, "f'_0"] \arrow[dr, phantom, "\text{cart}"] & X' \arrow[r] \arrow[d, "f'"] \arrow[dr, phantom, "\text{cart}"] & X \arrow[d, "f"] \\
Y'_1 \arrow[r, "p"'] & Y'_0 \arrow[r] & Y' \arrow[r] & Y & .
\end{tikzcd}
\]
Since $p$ is a $\DD$-equivalence, this indeed follows from a).

\indent\textit{b}) $\Rightarrow$ \textit{a}): As condition \textit{b}) is obviously stable under base change, we are reduced to proving that if $f$ satisfies \textit{b}), and if we have
\[
\begin{tikzcd}[column sep=large, row sep=large]
X \arrow[d, "f"'] \arrow[dr, phantom, "\text{cart}"] & X' \arrow[l, "q"'] \arrow[d, "f'"] \\
Y & Y' \arrow[l, "p"]
\end{tikzcd}
\tag*{$(*)$}
\]
if $p$ is a $\DD$-equivalence, then $q$ too, \textit{i.e.}
\[
\begin{tikzcd}[row sep=small]
H^*_{\DD}(X,M) \arrow[r, "\sim"] & H^*_{\DD}(X',M) \\
\wr| & \wr| \\
H^*_{\DD}(Y, f_*(M_X)) & H^*_{\DD}(Y', f'_*(M_{X'}))
\end{tikzcd}
\]
for every $M$ in $\mathcal{A} = \DD(e)$. But $F = f_*(M_X)$ is locally constant and is computed ``fiber by fiber'', likewise for $F' = f'_*(M_{X'})$.

\msp{page 81}

The equivalence \textit{a}) $\Leftrightarrow$ \textit{b}$'$) follows from \textit{a}) $\Leftrightarrow$ \textit{b}) by duality (condition \textit{a}) being self-dual), or by repeating the same argument for \textit{a}) $\Rightarrow$ \textit{b}$'$), and for \textit{b}$'$) $\Rightarrow$ \textit{a}) too (replacing cohomology by homology).

\textbf{Definition 2.} If the equivalent conditions of proposition 2 are satisfied, we say that $f$ is a $\DD$-\textit{fibration}.

N.B. This is a notion \textit{both} local and ``colocal'' on $Y$ (whereas left fibration is local, right fibration colocal).

\printedpage{50}
% head: A. Grothendieck : Dérivateurs, Ch. I, Édition des Universités de Montpellier II et Paris VII
% foot: Version 8 juillet 2001

\textbf{Proposition 3.} a) \textit{The notion of $\DD$-fibration is stable under base change and under composition.}\footnote[38]{Stability under composition \textit{is} not clear.}

\indent b) \textit{Let}
\[
\begin{tikzcd}[column sep=normal, row sep=large]
X \arrow[rr, "f"] \arrow[dr, "p"'] & \phantom{X} & X' \arrow[dl, "p'"] \\
& Y &
\arrow[from=1-2, to=2-2, phantom, "\text{\footnotesize comm}"]
\end{tikzcd}
\]
\textit{with $p,p'$ $\DD$-fibrations (it suffices that $p,p'$ be cohomological $\DD$-fibrations, or that they be homological $\DD$-fibrations, or that $p,p'$ be proper, or $p,p'$ smooth). If $f$ induces $\DD$-equivalences in the fibers, then $f$ is a $\DD$-equivalence.}\footnote[39]{Can be made considerably more precise in the case where $p$, $p'$ are $\DD$-fibrations, \textit{cf.} prop. 6, page 84.}

\indent\textsc{Proof.} \textit{a}) is immediate under condition \textit{a}) (prop. 2) of the notion of $\DD$-fibration. For \textit{b}), suppose that $f$ induces $\DD$-equivalences on the fibers. We want to prove that % [sic] On veux -> On veut, verb agreement
\[
H^*_{\DD}(X,M) \longrightarrow H^*_{\DD}(X',M)
\]
is an isomorphism, for every $M$ in $\mathcal{A}$, but this arrow is identified with
\[
H^*_{\DD}(Y, p_*(M_X)) \longrightarrow H^*_{\DD}(Y, p'_*(M_{X'}))
\]
deduced from
\[
p_*(M_X) \longrightarrow p'_*(M_{X'}) \qquad\qquad (\text{NB}\ \ M_X \simeq f^*M_{X'}\,)
\]
\msp{page 82}
and it suffices to prove that this arrow is an isomorphism. For this, we look fiber by fiber, and we use the fact that ``direct images are computed fiber by fiber''.

\textbf{Question.} Let $f : X \to Y$ be a $\DD$-fibration, and let $F \in \DD(Y)$. Then is $f_*f^*F$ computed by passing to the fibers? Same question for $f_!f^*F$. For the first property, concerning $f_*$, we show that it works if $F$ is of the form $p_*(M_Y)$, where $p : Y' \to Y$ is $\DD$-\textit{proper,} and $M \in \Ob\, \mathcal{A}$. On the other hand, the full subcategory $C$ of $\DD(Y)$ consisting of the $F$ for which it works is (unless I am mistaken) ``stable under $\DD$-$\mathsf{hot}$-$\varprojlim$'' in the following sense: if we have a category $I$ with final object $e$ forming a closed subcategory (\textit{i.e.} $e \to x \Rightarrow x = e$), and if we have $\mathcal{F} \in \DD(Y \times I)$, whence $I^\circ \to \DD(Y)$, $i \mapsto \mathcal{F}_i$ (skeletal diagram), if the $\mathcal{F}_i$ for $i \neq e$ are in $C$, and if $\mathcal{F}$ is in the essential image of $j_*$, where $j : Y \times (I \setminus \{e\}) \hookrightarrow Y \times I$ is the inclusion, then $\mathcal{F}_e$ is also in $C$.

\indent We may however consider that if we want $f : X \to Y$ to satisfy all the cohomological and homological properties that attach to the notion of a ``fibration'' (``locally trivial''), we should put among the conditions that $f$ is \textit{smooth}, and \textit{proper}, so that base changes by $f$ commute with the functors $g_*$, $g_!$. % [sic] on veux -> on veut, verb agreement

\printedpage{51}
% head: A. Grothendieck : Dérivateurs, Ch. I, Édition des Universités de Montpellier II et Paris VII
% foot: Version 8 juillet 2001

\msp{page 83}

Let us point out:

\textbf{Proposition 4.} \textit{Let $f : X \to Y$ be} proper. \textit{For $f$ to be a $\DD$-fibration, it is necessary and sufficient that for every $y_1 \to y_0$ in $Y$, the induced functor $X/y_1 \to X/y_0$ be a $\DD$-equivalence} (i.e. \textit{that $f$ be a} cohomological \textit{$\DD$-fibration}). \textit{Dually, if $f$ is $\DD$-smooth, for it to be a $\DD$-fibration, it is necessary and sufficient that for every $y_1 \to y_0$ in $Y$, $y_0\backslash X \to y_1\backslash X$ be a $\DD$-equivalence.}

\indent One could call \textit{perfect} $\DD$-fibrations the functors both proper and smooth that are cohomological $\DD$-fibrations (or what amounts to the same thing homological $\DD$-fibrations).

\textbf{Corollary.} \textit{Suppose $f$ a $\Cat$-cofibration (hence proper). For $f$ to be a $\DD$-fibration, it is necessary and sufficient that the functors $X_{y_0} \to X_{y_1}$ (for $y_0 \to u_1$ in $Y$) be $\DD$-equivalences. Dually, for $\Cat$-fibration (hence smooth).} % [sic] u_1 -> y_1, the arrow is y_0 -> y_1 ; [sic] Duallement -> Dualement, spelling

\textbf{Proposition 5.} \textit{Let $f : X \to Y$ be in} $\Dia$. \textit{Equivalent conditions:}
\begin{enumerate}[label={}, leftmargin=2em, labelsep=0pt]
\item a) \textit{$f$ is a $\DD$-equivalence, and remains so under every base change;}
\item b) \textit{$f$ is a $\DD$-fibration with $\DD$-contractible fibers.}
\end{enumerate}

\textsc{Proof.} \textit{a}) $\Rightarrow$ \textit{b}): That \textit{a}) implies that $f$ is a $\DD$-fibration is immediate from its definition (condition \textit{a}) of proposition 2), and that the fibers are contractible is tautological.

\textit{b}) $\Rightarrow$ \textit{a}): It suffices to show that \textit{b}) implies that $f$ is a $\DD$-equivalence. But this is a special case of proposition 3 \textit{b}) (taking $X' = Y$ there).

\msp{page 84}

When these conditions are satisfied, we expect that the functor
\[
f^* : \DD(Y) \longrightarrow \DD(X)
\]
is fully faithful, \textit{i.e.} that the adjunction morphisms % [sic] les morphisme -> les morphismes, plural agreement
\[
F \longrightarrow f_*f^*F \quad , \qquad f_!f^* \longrightarrow F % [sic] f_!f^* -> f_!f^*F, the object F is missing
\]
in $\DD(Y)$ are isomorphisms. (This is true if $f$ is $\DD$-proper or $\DD$-smooth, cf. cor. 3, p.71.) This again amounts to saying that the $f_*f^*F$ are computed by passing to the fibers, or again likewise for the $f_!f^*F$. This brings us back to the question on page 82.

\textbf{Definition 3.} We then say that $f$ is a \textit{trivial} $\DD$-fibration.

\textbf{Proposition 6} (makes prop. 3 page 81 more precise). \textit{Let us consider a commutative triangle in} $\Dia$
\[
\begin{tikzcd}[column sep=small, row sep=large]
X \arrow[rr, "f"] \arrow[dr, "p"'] & & X' \arrow[dl, "p'"] \\
& S & & & .
\end{tikzcd}
\]

\printedpage{52}
% head: A. Grothendieck : Dérivateurs, Ch. I, Édition des Universités de Montpellier II et Paris VII
% foot: Version 8 juillet 2001

\indent a) \textit{Suppose that $p$ and $p'$ are left $\DD$-fibrations, or that $p$ and $p'$ are right $\DD$-fibrations. Let $S_0$ be a subset of $S$ that meets every connected component of $S$. For $f$ to be a $\DD$-equivalence, it} suffices \textit{that for every $s_0 \in S_0$, $f$ induce a $\DD$-equivalence $X_{s_0} \to Y_{s_0}$.}

\indent b) \textit{This condition is also necessary}\footnote[40]{Not proved!\footnotemark[41] Even if $X' = S$, it is not clear, even if $\DD = \mathsf{HOTAB}$, $f$ proper and smooth. But OK if $S$ is $\DD$-contractible.}\footnotetext[41]{\underline{False}, \textit{cf.} counterexample, p. 87.} \textit{if $p$ and $p'$ are $\DD$-fibrations. So in this case, $f$ is a $\DD$-equivalence if and only if $f$ induces $\DD$-equivalences fiber by fiber.}

\msp{page 85}

\textsc{Proof.} a) It suffices, by virtue of proposition 3, to show that the hypothesis implies that for \textit{every} $s \in S$, $f$ induces a $\DD$-equivalence $f_s : X_s \to X'_s$. This follows at once from the lemma:

\textbf{Lemma.} \textit{If $p$ and $p'$ are both either left $\DD$-fibrations or right $\DD$-fibrations, and if $s_1 \to s_0$ is an arrow in $S$, then $f_{s_1}$ is a $\DD$-equivalence if and only if $f_{s_0}$ is. (Thus the set $S' \subset \Ob\, S$ of $s \in S$ such that $f_s$ is a $\DD$-equivalence is a union of connected components of $S$.)}

\indent Suppose for example that $p,p'$ are left $\DD$-fibrations. Let us consider, for $s \in S$, the diagram
\[
\begin{tikzcd}[row sep=large]
X_s \arrow[r, "f_s"] \arrow[d] & X'_s \arrow[d] \\
X/s \arrow[r, "f/s"'] & X'/s & .
\end{tikzcd}
\]
The vertical arrows are $\DD$-equivalences by the hypothesis on $p,p'$, so $f_s$ is a $\DD$-equivalence if and only if $f/s$ is. Let us now consider
\[
\begin{tikzcd}[row sep=large]
X/s_1 \arrow[r, "f/s_1"] \arrow[d] & X'/s_1 \arrow[d] \\
X/s_0 \arrow[r, "f/s_0"'] & X'/s_0 & .
\end{tikzcd}
\]
Here again, the hypothesis on $p,p'$ implies that the vertical arrows are $\DD$-equivalences, so $f_{s_1}$ is a $\DD$-equivalence if and only if $f_{s_0}$ is.

\msp{page 86}

If $Y' = S$, the statement (not proved) says that if $f : X \to S$ is a $\DD$-fibration and a $\DD$-equivalence, then the fibers are $\DD$-contractible. This is not certain even for a projection $X = T \times S \to S$. If $\eta = \mathit{hot}_0(S)$, $\xi = \mathit{hot}_0(T)$, the question reduces to this one: If for every $M$ in $\mathcal{A}$, we have: the homomorphism
\[
\alpha_M : \xi \times_{\DD} M \longrightarrow M
\]
\printedpage{53}
% head: A. Grothendieck : Dérivateurs, Ch. I, Édition des Universités de Montpellier II et Paris VII
% foot: Version 8 juillet 2001
satisfies
\[
\eta \times_{\DD} \alpha_M : \eta \times_{\DD} (\xi \times_{\DD} M) \xrightarrow{\sim} \eta \times_{\DD} M
\]
is an isomorphism, then $\alpha_M$ is itself an isomorphism.

\msp{page 87}

Here is a \textit{counterexample}\footnote[42]{See an even simpler counterexample on pages 93, 94. (But it has the drawback of not having connected fibers.) % [sic] inconvenient -> inconvénient, missing accent
} which I first give in the context of topological spaces, instead of $\Cat$. I am going to construct a locally trivial and nontrivial fibration
\[
\begin{tikzcd}[column sep=0.2em, row sep=large]
X \arrow[d, "f"'] & \supset S^1 & \\
S & = S^1 & \qquad ,
\end{tikzcd}
\]
with fiber $S^1$, such that $f$ is a $\DD$-equivalence for every $\DD = \mathsf{Der}(k\text{-}\mathsf{Mod})$, $k$ being a ring in which 2 is invertible. This means that
\[
H_*(X,k) \longrightarrow H_*(S,k)
\]
is an isomorphism modulo 2-torsion (kernel and cokernel are annihilated by a power of 2 - in fact by 2 here), or again $H_*(X,k) \xrightarrow{\sim} H_*(S,k)$ if 2 invertible in $k$. But the fiber $S^1$ is not $\DD$-contractible (unless $k = 0$!).\footnote[43]{Much simpler example still: we take $X = S^1$, and $X \to S$ the two-sheeted covering. If we take the $n$-sheeted covering, it is also OK if $k$ a ring in which $n$ is invertible.}

\indent The fibration is given up to unique isomorphism, by fixing a point $s \in S$, and giving oneself an automorphism $u$ of the fiber $X_s$ (``monodromy''). (The isomorphism class of the bundle depends only on the isotopy class of $u$). I take for $u$ the symmetry (reversing orientation) $z \mapsto z^{-1}$ of the standard circle $\{z \in \mathbb{C} \mid |z| = 1\}$. This fibration ``is realized'' by ordered sets $X \to S$, $S^1$ being represented (on the base, as on the fibers) by % [sic] Je prend -> Je prends, verb agreement ; [sic] ordonés -> ordonnés, spelling
\[
\begin{tikzcd}[column sep=normal, row sep=normal]
& & a & & \\
S \ = & 0 \arrow[ur] \arrow[dr] & & 0' \arrow[ul] \arrow[dl] & = \ S' \amalg_R S'' \quad , \\
& & b & &
\end{tikzcd}
\]
where $R$ is the discrete category having two elements $a$, $b$. We define $X$ as $X' \amalg_Y X''$, where $X' = X|S'$, $X'' = X|S''$, $Y = X|R$.\footnote[44]{N.B. $S'$, $S''$, $R = S' \cap S''$ closed in $S$, $X'$, $X''$, $Y = X' \cap X''$ closed in $X$.}